\textbf{To: Show Management}\\
\textbf{Subject: Recommendation for a More Professional Yet Watchable Voting System}

Our analyses indicate that controversy risk concentrates in weeks where judges' evaluations and audience preferences diverge.
To prioritize technical merit while preserving an audience role, we recommend piloting \textit{Professional Bottom-2} (\textit{PB2}), a judges-gated two-stage voting policy,
validated by replay-based stress tests: (i) pairwise Jaccard similarity across Percent/Rank/B2/PB2, (ii) controversial-case replays, and (iii) a unified bias test.

\paragraph{Baseline reference (current mechanism).}
The incumbent format can be summarized as ``combined Bottom-2 + judges' save'': contestants are first ranked by a combined
judge--fan score, the Bottom-2 form an at-risk set, and judges decide which of the two is eliminated.

\paragraph{Recommended policy (PB2: \textit{Professional Bottom-2}): judges-only gate $\rightarrow$ fan-aware decision within the gate.}
Under PB2, each week proceeds in two stages:
\begin{enumerate}
  \item \textbf{Judges-only gate (risk set).} Rank contestants by judges' performance share $q^J_{i,t}$ and designate the Bottom-2
  (or Bottom-$k$ in late-stage variants) as the at-risk set.
  \item \textbf{Fan-aware decision within the gate.} Within the at-risk set, eliminate the contestant with the lower transparent blended score
  \[
  S_{i,t}=w_J\, q^J_{i,t} + (1-w_J)\, p_{i,t},
  \]
  where $p_{i,t}$ denotes the inferred/proxy fan vote share.
\end{enumerate}

\paragraph{Why this improves professionalism.}
Because elimination is always decided among the judges-defined bottom performers, extreme popularity cannot keep a weak performer
out of the at-risk set. This makes professionalism a structural property of the system rather than a tunable side effect.

\paragraph{How we preserve watchability (without randomness).}
To maintain suspense while remaining performance-grounded, we recommend a phase-dependent schedule:
early weeks use Bottom-2 gating to emphasize professionalism; late in the season (when remaining contestants are few and performance gaps narrow),
management may enable structured suspense by expanding the gate to a Bottom-$k$ band and/or using a small number of structured multi-elimination weeks,
with eliminations restricted to the judges-defined low-performance band.

\paragraph{Management knobs and deployment.}
PB2 offers interpretable control knobs:
(i) $w_J$ (judge--fan balance within the gate),
(ii) the late-stage trigger (e.g., remaining contestants $\le K$ or week threshold $t_0$),
and (iii) the late-stage gate size $k$ (and optional multi-elimination placement).
We recommend selecting a robust operating point supported by replay-based stress tests (Jaccard similarity, controversial-case replays, and a unified bias test), and deploying PB2 in a one-season
\emph{shadow mode} computed in parallel with the broadcast rule. This yields an auditable comparison of eliminations, finalists, and winner outcomes
before any on-air adoption, supported by a clear public explanation of the two-stage logic and the role of audience influence.
